%%
%% Test: Color System — v1.1
%% Stage: colors refactor (stages 1-4)
%% 
%% Purpose: Verify color system:
%%   - All matin* colors are defined (CMYK)
%%   - Colors are defined BEFORE consumers (order check)
%%   - Colors work in text, boxes, headings, graphics
%%   - Dominant color: matinblue
%%   - CMYK color model is used
%% 

\documentclass{matinbook}

\title{تست سیستم رنگ — نسخه ۱.۱}
\author{تیم MatinBook}
\date{۱۴۰۵}

\booktitle{تست رنگ‌ها}

\begin{document}

\maketitle

\chapter{بررسی سیستم رنگ}
\label{chap:test-colors}

این سند، سیستم رنگ \texttt{mb-theme-colors} را بررسی می‌کند.

%----------------------------------------
\section{رنگ‌های اصلی}
\label{sec:primary-colors}

رنگ‌های اصلی MatinBook:

\begin{itemize}
    \item \textcolor{matinblue}{\textbf{matinblue}} — رنگ غالب (آبی)
    \item \textcolor{matindarkblue}{\textbf{matindarkblue}} — آبی تیره
    \item \textcolor{matinlightblue}{\textbf{matinlightblue}} — آبی روشن
\end{itemize}

%----------------------------------------
\section{رنگ‌های ثانویه}
\label{sec:secondary-colors}

\begin{itemize}
    \item \textcolor{matingreen}{\textbf{matingreen}} — سبز (تعاریف)
    \item \textcolor{matindarkgreen}{\textbf{matindarkgreen}} — سبز تیره
    \item \textcolor{matinlightgreen}{\textbf{matinlightgreen}} — سبز روشن
\end{itemize}

%----------------------------------------
\section{رنگ‌های تأکیدی}
\label{sec:accent-colors}

\begin{itemize}
    \item \textcolor{matinorange}{\textbf{matinorange}} — نارنجی (مثال‌ها)
    \item \textcolor{matindarkorange}{\textbf{matindarkorange}} — نارنجی تیره
    \item \textcolor{matinlightorange}{\textbf{matinlightorange}} — نارنجی روشن
    \item \textcolor{matinred}{\textbf{matinred}} — قرمز (هشدارها)
    \item \textcolor{matindarkred}{\textbf{matindarkred}} — قرمز تیره
    \item \textcolor{matinlightred}{\textbf{matinlightred}} — قرمز روشن
    \item \textcolor{matinpurple}{\textbf{matinpurple}} — بنفش (تمرین‌ها)
    \item \textcolor{matindarkpurple}{\textbf{matindarkpurple}} — بنفش تیره
    \item \textcolor{matinlightpurple}{\textbf{matinlightpurple}} — بنفش روشن
\end{itemize}

%----------------------------------------
\section{رنگ‌های خنثی}
\label{sec:neutral-colors}

\begin{itemize}
    \item \textcolor{matindarkgray}{\textbf{matindarkgray}} — خاکستری تیره
    \item \textcolor{matinmediumgray}{\textbf{matinmediumgray}} — خاکستری متوسط
    \item \textcolor{matinlightgray}{\textbf{matinlightgray}} — خاکستری روشن
    \item \textcolor{matinbordergray}{\textbf{matinbordergray}} — خاکستری حاشیه
    \item \textcolor{matincream}{\textbf{matincream}} — کرم
\end{itemize}

%----------------------------------------
\section{رنگ در جعبه‌ها}
\label{sec:colors-in-boxes}

\begin{theorem}[قضیه نمونه]
    برای هر عدد حقیقی $x$، داریم $x^2 \geq 0$.
\end{theorem}

\begin{definition}[عدد اول]
    عدد طبیعی بزرگ‌تر از ۱ که تنها دو مقسوم‌علیه مثبت دارد.
\end{definition}

\begin{example}[مثال نمونه]
    اعداد $2, 3, 5, 7, 11$ نمونه‌هایی از اعداد اول هستند.
\end{example}

\begin{remark}[نکته]
    عدد ۱ نه اول است و نه مرکب.
\end{remark}

%----------------------------------------
\section{رنگ در ریاضیات}
\label{sec:colors-in-math}

\begin{equation}
    \label{eq:colors-test}
    \textcolor{matinblue}{f(x)} = \textcolor{matingreen}{x^2} + \textcolor{matinorange}{2x}
\end{equation}

%----------------------------------------
\section{رنگ در جدول}
\label{sec:colors-in-table}

\begin{table}[h]
\centering
\caption{جدول رنگی}
\label{tab:colors-test}
\begin{tabular}{L{3cm} C{2cm} R{3cm}}
\toprule
\rowcolor{matinlightblue}
\textbf{ستون ۱} & \textbf{ستون ۲} & \textbf{ستون ۳} \\
\midrule
\rowcolor{matinlightgray}
متن & ۱۰ & ۲۰ \\
\rowcolor{matincream}
متن & ۲۰ & ۳۰ \\
\bottomrule
\end{tabular}
\end{table}

%----------------------------------------
\section{رنگ در گرافیک}
\label{sec:colors-in-graphics}

\begin{latin}
\begin{tikzpicture}
    \fill[matinblue] (0,0) circle (0.5cm);
    \fill[matingreen] (1.5,0) circle (0.5cm);
    \fill[matinorange] (3,0) circle (0.5cm);
    \fill[matinred] (4.5,0) circle (0.5cm);
    \fill[matinpurple] (6,0) circle (0.5cm);
\end{tikzpicture}
\end{latin}

%----------------------------------------
\section{نتیجه‌گیری}
\label{sec:conclusion}

اگر این سند بدون خطا کامپایل شود:

\begin{itemize}
    \item ✅ همه‌ی ۲۰ رنگ \texttt{matin*} تعریف شده‌اند
    \item ✅ رنگ‌ها در CMYK تعریف شده‌اند
    \item ✅ \texttt{matinblue} رنگ غالب است
    \item ✅ رنگ‌ها در جعبه‌ها کار می‌کنند
    \item ✅ رنگ‌ها در ریاضیات کار می‌کنند
    \item ✅ رنگ‌ها در جدول کار می‌کنند
    \item ✅ رنگ‌ها در گرافیک کار می‌کنند
    \item ✅ ترتیب بارگذاری رنگ‌ها درست است
\end{itemize}

\textbf{وضعیت تست:} قبول

\end{document}
